\subsection{固定した低 row と column からの三高分母の完全復元}
\label{nc17:fixed-low}

今後は三高分母 core の全分類を待たず，隣接 core から低データが決まった時点で
次の二次復元を使える．低 ray を $0$ とし，整数
\[
 r_i=p_{0i},\quad x_i=p_{i0},\quad
 \sum_i r_i=-1,\quad q_i=1+r_ix_i>1\quad(i=1,2,3)
\]
を固定する．三つの $r_i$ の同符号の対を $r_2,r_3$ とし，
$s_i=-q_i-x_i<0$，$Q_{ij}=q_iq_j-1>0$ と置く．
$A=p_{12}$ は
\[
 (s_2A-Q_{12})\{s_3(s_1-A)-Q_{13}\}
 -Q_{23}A(s_1-A)=0
\]
を満たす．二次係数 $a=Q_{23}-s_2s_3$ は，$r_2,r_3>0$ なら負，
$r_2,r_3<0$ なら正なので消えない．従って高高 entries の completion は高々二つである．
その整数根から
\[
 p_{21}=Q_{12}/A,\quad p_{13}=s_1-A,\quad p_{31}=Q_{13}/(s_1-A),
\]
\[
 p_{23}=s_2-Q_{12}/A,\qquad p_{32}=s_3-Q_{13}/(s_1-A)
\]
を復元する．必要十分な整数条件は $A\ne0,s_1$，$A\mid Q_{12}$，
$s_1-A\mid Q_{13}$ であり，その後に全正 cycle の条件を課す．
判別式は
\[
 \Delta=(s_1s_2s_3-s_1Q_{23}-s_2Q_{13}-s_3Q_{12})^2
            -4Q_{12}Q_{13}Q_{23}.
\]
任意の探索上限を置かず，整数平方根と二つの割り算だけで全候補を得る．

実装 \nolinkurl{research/exact_three_high_completion_v17.py} の
\texttt{complete\_three\_high} は元の物理標識の順を保存する．
同符号対を選ぶ内部の標識交換を，隣接 core への clutching に誤って持ち込まない．
通常の整数行列だけを求める場合と，正の weights を求める場合を区別できる．
\nolinkurl{verify_three_high_completion_v17.py} は別の signed-divisor 復元との一致も検査する．

分母 $q$ と低 row $r$ が既知なら，$x_i=(q_i-1)/r_i$ も決まるため同じ手順を使える．
一方，分母を可変のまま低 row だけを固定した場合には高々二つとはいえない．
既存の rank 二 Schur 無限族は低 column 自体が変化する例である．
$r=(-2,-2,3)$，$x=(-3,-3,2)$，$q=(1,7,7,7)$ は二つの正の標識付きcompletionを持つが，
両者は高 ray $1,2$ の交換で対応する．標識付き解の個数を幾何学的同値類数と同一視しない．

% Reusable finite bound; optional research-note material, not needed for K4 exclusion.
\subsection{反対符号の二つの例外 reverse が大きい場合の有限化}
\label{fc17:tail-note}

低 row と column が
\[
 r=(P,-P+h-1,-h),\qquad x=(1,-1,-v),\qquad
 P\ge h\ge1,\quad v\ge1
\]
である整数三高分母模型を考える．分母は
\[
 q=(P+1,P-h+2,hv+1)
\]
となる．高高 entry $t=p_{13}$ を用いて他の entries を行和・pair式から消去すると，
\[
 P\{vt^2+2(hv+1)t+(h-1)(hv+1)\}
 +vt^2+\{v(h+1)+1-h\}t+hv(h-1)=0
\]
を得る．$t\ne0$ であり，$t>0$ なら左辺は正なので，$t=-z$，$z\ge1$ でなければならない．
\[
 \begin{split}
 A(z)&=vz^2-2(hv+1)z+(h-1)(hv+1),\\
 B(z)&=vz^2-\{v(h+1)+1-h\}z+hv(h-1)
 \end{split}
\]
と置けば $PA(z)+B(z)=0$ であり，
\[
 B(z)-A(z)=\{(h-1)v+h+1\}z-(h-1)>0.
\]
従って $A(z)<0$ で，
\[
 P+1=\frac{\{(h-1)v+h+1\}z-(h-1)}{-A(z)}.
\]
$A(z)=B(z)=0$ という共鳴は生じない．

$z\ge2h+2$ なら
\[
 A(2h+2)=(v-1)(h^2+3h+4)+(h^2-1)\ge0
\]
であり，その先で $A$ は増加する．また
\[
 B(z)=vz(z-h-1)+(h-1)z+hv(h-1)>0
\]
なので矛盾する．従って
\[
 1\le z\le2h+1,\qquad
 P\le\{(h-1)v+h+1\}(2h+1)-h.
\]
よって $h,v$ を固定すれば全 $P$ を完全に復元できる．
全分母が奇数なら $P$ は偶数，$h$ は奇数である．
通常 reverse の上界 $h\le4$ が使える場合には $h=1,3$ のみとなる．
この有限化は他の graph で反対符号の大きい例外 entries が残る場合に転用できる．

